\scanpage{138}
donc $H^0(X,\underline\omega^2)^{\mathrm{rég}}=0$ :
\uline{Il n'y a qu'une seule connexion projective} \inserted{régulière à l'$\infty$} \uline{sur $X$} si
\uline{$X$ est de genre $0$ et à $\nu\leq3$ trous}. Et
par contre $\nu\geq4$, l'espace $H^0(X,\underline\omega^2)^{\mathrm{rég}}$ est
de dim $\nu-3$ -- c'est encore la dim de
l'espace modulaire pour la droite projective
à $\nu$ trous, et sans doute on trouve
encore un isom can $H^0(X,\underline\omega^2)^{\mathrm{rég}}
\simeq\Omega^1_{M_{0,\nu}}(s)$ \ldots\ Bien entendu, dans le
cas $\nu\leq3$, la seule et unique connexion
projective sur $X$ est celle induite par
l'unique conn. projective de $\widehat X$ (et leur
description en \unsure{termes|tant} de conn. n'a rien
de transcendant !). Dans le cas de $\nu\leq3$,
% slanted
\margin{\unsure{Les|ses} \unsure{résultats} restent bien \unsure{entendu} \unsure{vrais} sur une base $S$ \unsure{quelconque} \unsure{au lieu de} \unsure{$\mathbb C$}.}
\deleted{\unsure{les}} la connexion projective \inserted{unique} \deleted{est d'ailleurs}
$\mathbb P_X=\widehat X\times X$, muni de la section diagonale \inserted{\unsure{(évidente)|(d'ailleurs)}} et
de la connexion triviale. Dans le cas
de $X$ quelconque, \inserted{(de genre $0$)} cette
connexion projective est tjrs définie, et
fournit une \uline{section canonique}.
% archive sheet number: 137
